% ============================================================
%  OOD-gated TTA MoE 방법 정리 슬라이드 초고
%  컴파일: xelatex ood_tta_draft.tex
% ============================================================
\documentclass[aspectratio=169]{beamer}
\usetheme{Madrid}
\usecolortheme{seahorse}

\usepackage{booktabs}
\usepackage{amsmath}
\usepackage{amssymb}
\usepackage{xcolor}
\usepackage{array}
\usepackage{tikz}
\usepackage{pgfplots}
\pgfplotsset{compat=1.18}
\usetikzlibrary{arrows.meta, positioning, shapes.geometric, fit, calc, backgrounds}
\usepackage{fontspec}
\usepackage{xeCJK}
\setCJKmainfont[AutoFakeSlant=0.2, AutoFakeBold=2]{Noto Sans CJK KR}
\setCJKsansfont[AutoFakeSlant=0.2, AutoFakeBold=2]{Noto Sans CJK KR}
\setCJKmonofont{Noto Sans CJK KR}

\definecolor{mainblue}{RGB}{55,105,180}
\definecolor{maingreen}{RGB}{45,145,85}
\definecolor{mainorange}{RGB}{205,125,30}
\definecolor{mainred}{RGB}{185,55,55}
\definecolor{lightgray}{RGB}{245,245,245}

\setbeamertemplate{navigation symbols}{}
\setbeamertemplate{itemize items}[circle]
\setbeamerfont{footnote}{size=\tiny}

\title{OOD-aware Stability-gated Expert MoE}
\subtitle{CIC-IDS2017 불균형 공격 분류를 위한 Expert Gate, 외부 OOD 학습, 진단 체계}
\author{westpark}
\date{2026-06-17}
\institute{CIC-IDS2017 \quad / \quad Auxiliary OOD: CIC2018 후보}

\begin{document}

\begin{frame}
  \titlepage
\end{frame}

\begin{frame}{문제의식}
\begin{itemize}
  \item 기존 MoE/TTA 실험은 rare attack recall 개선 가능성을 보였지만, expert 선택 실패 원인 해석이 어려웠다.
  \item 새 구조에서는 benign을 expert에 배정하지 않고, 공격 클래스만 disjoint expert로 분할한다.
  \item 핵심 질문은 세 가지다.
\end{itemize}

\vspace{0.4em}
\begin{block}{진단해야 할 병목}
\begin{enumerate}
  \item expert가 자기 공격 샘플을 OOD로 reject하는가?
  \item OOD 또는 benign이 특정 expert로 잘못 accept되는가?
  \item expert가 선택된 뒤 local classifier가 내부 공격 클래스를 구분하지 못하는가?
\end{enumerate}
\end{block}
\end{frame}

\begin{frame}{제안 방법: 전체 구조}
\begin{center}
\resizebox{0.96\textwidth}{!}{\begin{tikzpicture}[
  font=\small,
  node distance=0.7cm and 0.9cm,
  data/.style={draw, rounded corners=3pt, fill=blue!8, minimum width=2.5cm, minimum height=0.72cm, align=center},
  proc/.style={draw, rounded corners=3pt, fill=green!10, minimum width=2.7cm, minimum height=0.72cm, align=center},
  gate/.style={draw, rounded corners=3pt, fill=orange!12, minimum width=2.7cm, minimum height=0.72cm, align=center},
  eval/.style={draw, rounded corners=3pt, fill=red!8, minimum width=2.7cm, minimum height=0.72cm, align=center},
  expert/.style={draw, rounded corners=3pt, fill=gray!10, minimum width=2.55cm, minimum height=0.62cm, align=center},
  arr/.style={-{Stealth[length=2.2mm]}, thick},
  dashedarr/.style={-{Stealth[length=2.2mm]}, thick, dashed},
  lab/.style={font=\scriptsize\bfseries, align=center}
]

\node[data] (id) {ID 데이터\\CIC-IDS2017};
\node[proc, right=of id] (split) {Train / Val / Test\\stratified split};
\node[proc, right=of split] (partition) {공격 클래스 분할\\Benign 제외};

\node[data, below=1.0cm of id] (aux) {Auxiliary OOD\\예: CIC2018 train};
\node[data, below=1.0cm of split] (sc) {SC-OOD 평가셋\\예: CIC2018 test};

\node[expert, right=1.1cm of partition, yshift=0.8cm] (e1) {Expert 1\\DoS / DDoS};
\node[expert, below=0.18cm of e1] (e2) {Expert 2\\Scan / Brute};
\node[expert, below=0.18cm of e2] (e3) {Expert 3\\Web / Bot / Rare};

\node[gate, right=1.1cm of e2] (local) {Local XGBoost\\owned attack 분류};
\node[gate, above=0.55cm of local] (oodgate) {OE Gate\\ID vs OOD};
\node[gate, below=0.55cm of local] (sweep) {Expert별 threshold\\validation sweep};

\node[eval, right=1.1cm of local] (predict) {Test-time decision\\accept $\rightarrow$ select};
\node[eval, above=0.52cm of predict] (benign) {all reject\\$\Rightarrow$ Benign};
\node[eval, below=0.52cm of predict] (diag) {진단 CSV\\병목 / OOD accept};

\draw[arr] (id) -- (split);
\draw[arr] (split) -- (partition);
\draw[arr] (partition) -- (e1.west);
\draw[arr] (partition) -- (e2.west);
\draw[arr] (partition) -- (e3.west);

\draw[dashedarr] (aux) -| node[pos=0.23, left, font=\scriptsize] {negative} (oodgate.south west);
\draw[dashedarr] (sc) -| node[pos=0.22, right, font=\scriptsize] {평가 전용} (diag.south west);

\draw[arr] (e1.east) -- (oodgate.west);
\draw[arr] (e2.east) -- (local.west);
\draw[arr] (e3.east) -- (sweep.west);
\draw[arr] (oodgate) -- (predict);
\draw[arr] (local) -- (predict);
\draw[arr] (sweep) -- (predict);
\draw[arr] (predict) -- (benign);
\draw[arr] (predict) -- (diag);

\begin{scope}[on background layer]
  \node[draw=green!45!black, rounded corners=4pt, fit=(e1)(e2)(e3), inner sep=0.18cm,
        label={[lab, text=green!35!black]above:Disjoint attack experts}] {};
  \node[draw=orange!70!black, rounded corners=4pt, fit=(oodgate)(local)(sweep), inner sep=0.22cm,
        label={[lab, text=orange!60!black]below:gate calibration}] {};
\end{scope}

\node[lab, above=0.15cm of id, text=blue!55!black] {ID 학습/평가};
\node[lab, left=0.12cm of aux, text=orange!70!black] {외부 OOD};

\end{tikzpicture}
}
\end{center}
\end{frame}

\begin{frame}{모델 메커니즘}
\begin{columns}[T]
\column{0.48\textwidth}
\textbf{학습 단계}
\begin{itemize}
  \item ID 데이터: CIC-IDS2017
  \item Benign은 expert에 배정하지 않음
  \item 공격 클래스만 taxonomy 기반으로 분할
  \item 각 expert는 owned attack class만 local classifier로 학습
  \item OE gate는 owned class vs negative를 binary로 학습
\end{itemize}

\column{0.48\textwidth}
\textbf{추론 단계}
\begin{itemize}
  \item 각 expert가 독립적으로 accept/reject 수행
  \item 모든 expert가 reject하면 benign 예측
  \item 하나만 accept하면 해당 expert 선택
  \item 여러 expert가 accept하면 gate margin과 TTA stability로 선택
  \item 선택된 expert의 local 예측을 global label로 변환
\end{itemize}
\end{columns}
\end{frame}

\begin{frame}{OOD 학습 재정의}
\begin{block}{기존 OE gate의 한계}
기존 negative는 같은 데이터셋 내부의 non-owning class였다. 따라서 엄밀한 open-set OOD라기보다 \textbf{owning class vs not}에 가까웠다.
\end{block}

\vspace{0.4em}
\begin{block}{수정된 방향}
ImOOD에서 TinyImages80M을 auxiliary OOD training data로 쓰는 위치에 대응하여, IDS 실험에서는 feature-aligned 외부 트래픽 데이터셋을 OE gate negative로 추가한다.
\end{block}

\vspace{0.4em}
\begin{center}
\begin{tabular}{lll}
\toprule
역할 & 이미지 OOD 예시 & IDS 실험 대응 \\
\midrule
ID train/test & CIFAR/ImageNet & CIC-IDS2017 \\
Auxiliary OOD train & TinyImages80M & CIC2018 train split \\
SC-OOD eval & SVHN/LSUN/Places & CIC2018 test split 또는 다른 IDS \\
\bottomrule
\end{tabular}
\end{center}
\end{frame}

\begin{frame}[fragile]{OOD 데이터 준비}
\textbf{새 준비 스크립트}
\begin{itemize}
  \item \texttt{data/prepare\_external\_ood.py}
  \item ID 데이터셋 feature schema를 기준으로 source OOD feature를 재정렬
  \item CICFlowMeter 2017/2018 컬럼명 차이를 alias rule로 보정
  \item CIC2017 기준 CIC2018 feature alignment: 70/70 matched 확인
\end{itemize}

\vspace{0.5em}
\textbf{생성 예시}
\begin{verbatim}
python src/prepare_external_ood.py \
  --id-data data/cic2017_proc.pkl \
  --source-data data/cic2018_proc.pkl \
  --out data/cic2018_ood_for_cic2017.pkl \
  --max-samples 500000
\end{verbatim}
\end{frame}

\begin{frame}{Expert별 Threshold Sweep}
\begin{itemize}
  \item 기존 방식: 모든 expert가 같은 \texttt{threshold\_quantile} 사용
  \item 문제: DoS/Scan은 낮은 threshold가 유리하지만, WebBot/Rare는 낮추면 precision이 크게 붕괴
  \item 수정: validation split에서 expert마다 threshold 후보를 sweep
\end{itemize}

\vspace{0.3em}
\begin{block}{목적함수 직관}
\[
  \text{score}
  = \text{owned accept}
  - \lambda_b \cdot \text{benign accept}
  - \lambda_o \cdot \text{non-owned attack accept}
\]
\end{block}

\vspace{0.3em}
\textbf{저장 파일}
\begin{itemize}
  \item \texttt{expert\_threshold\_sweep.csv}: 모든 후보 threshold 결과
  \item \texttt{expert\_threshold\_selected.csv}: expert별 최종 선택 threshold
\end{itemize}
\end{frame}

\begin{frame}{SC-OOD Accept-rate 평가}
\begin{itemize}
  \item \texttt{external\_ood\_data}: gate 학습에 들어가는 auxiliary OOD
  \item \texttt{sc\_ood\_data}: 학습 후 accept/reject 성능만 보는 평가용 OOD
  \item 같은 파일을 둘 다 쓰면 sanity check이고, 논문식 평가는 train/test 또는 dataset 분리가 필요하다.
\end{itemize}

\vspace{0.4em}
\begin{center}
\begin{tabular}{ll}
\toprule
파일 & 해석 \\
\midrule
\texttt{sc\_ood\_accept\_summary.csv} & 전체 OOD all-reject / any-accept \\
\texttt{sc\_ood\_expert\_accept.csv} & expert별 OOD accept, selected, margin 분포 \\
\bottomrule
\end{tabular}
\end{center}
\end{frame}

\begin{frame}{결과 파일 해석 체계}
\begin{center}
\begin{tabular}{lll}
\toprule
파일 & 주요 컬럼 & 용도 \\
\midrule
\texttt{run\_summary.csv} & macro-F1, all-reject, threshold & 실험 전체 요약 \\
\texttt{bottleneck\_by\_class.csv} & primary\_bottleneck & class별 병목 식별 \\
\texttt{expert\_summary.csv} & owned/benign/non-owned accept & expert 단위 gate 해석 \\
\texttt{activation\_matrix.csv} & class별 expert accept/selected & 상세 activation 확인 \\
\texttt{selection\_reason.csv} & all\_reject/only\_accept/tiebreak & TTA 선택 개입 여부 \\
\bottomrule
\end{tabular}
\end{center}

\vspace{0.6em}
\begin{block}{핵심 판독법}
\begin{itemize}
  \item \texttt{gate\_reject}: threshold가 너무 보수적
  \item \texttt{local\_classifier}: expert는 맞게 선택됐지만 내부 분류 실패
  \item \texttt{benign\_false\_accept}: threshold가 너무 느슨하거나 OOD gate 약함
\end{itemize}
\end{block}
\end{frame}

\begin{frame}[fragile]{실험 명령 예시}
\textbf{1단계: CIC2018 OOD 준비}
\begin{verbatim}
python src/prepare_external_ood.py \
  --id-data data/cic2017_proc.pkl \
  --source-data data/cic2018_proc.pkl \
  --out data/cic2018_ood_for_cic2017.pkl \
  --max-samples 500000
\end{verbatim}

\vspace{0.3em}
\textbf{2단계: Expert threshold sweep + auxiliary OOD 학습}
\begin{verbatim}
python src/code_ood_v1.py \
  --data data/cic2017_proc.pkl \
  --ood_gate oe \
  --threshold_mode expert_sweep \
  --external_ood_data data/cic2018_ood_train_for_cic2017.pkl \
  --sc_ood_data data/cic2018_ood_test_for_cic2017.pkl
\end{verbatim}
\end{frame}

\begin{frame}{현재까지의 결론}
\begin{itemize}
  \item 기존 q=0.05 결과의 주된 성능 하락은 TTA가 아니라 gate reject였다.
  \item q=0.01은 attack recall을 회복하지만, q=0.001은 WebBot/Rare precision 붕괴를 유발했다.
  \item 따라서 공통 threshold가 아니라 expert별 threshold가 필요하다.
  \item OOD 학습도 in-dataset non-owning class만으로는 부족하므로 external auxiliary OOD를 negative로 추가했다.
  \item 앞으로는 ID 성능, SC-OOD reject, class별 병목을 동시에 보고 개선해야 한다.
\end{itemize}
\end{frame}

\begin{frame}{다음 실험 체크리스트}
\begin{enumerate}
  \item CIC2018 OOD를 train/test artifact로 분리
  \item \texttt{expert\_sweep + external\_ood\_data} run 실행
  \item \texttt{run\_summary.csv}에서 macro-F1과 all-reject 확인
  \item \texttt{expert\_threshold\_selected.csv}에서 expert별 threshold 확인
  \item \texttt{sc\_ood\_accept\_summary.csv}에서 SC-OOD reject 확인
  \item WebBot/Rare가 OOD sink가 되는지 \texttt{sc\_ood\_expert\_accept.csv}로 확인
\end{enumerate}
\end{frame}

\end{document}
